% 图契约：本题的源、边界和推导逐项对应；箭头符号按图注解释。
\documentclass[tikz,border=7pt]{standalone}
\usepackage{fontspec,xeCJK,amsmath,amssymb}
\IfFontExistsTF{Times New Roman}{\setmainfont{Times New Roman}}{\setmainfont{TeX Gyre Termes}}
\IfFontExistsTF{Heiti SC}{\setCJKmainfont[AutoFakeBold=1.5]{Heiti SC}}{\setCJKmainfont[AutoFakeBold=1.5]{FandolHei-Regular}}
\usetikzlibrary{arrows.meta,bending,calc,patterns,positioning}
\definecolor{ink}{HTML}{243745}
\definecolor{blue}{HTML}{275DAB}
\definecolor{green}{HTML}{177C63}
\definecolor{red}{HTML}{BD4935}
\definecolor{orange}{HTML}{AA6C16}
\definecolor{purple}{HTML}{715393}
\tikzset{every node/.style={font=\fontsize{9.3}{12.5}\selectfont,text=ink},
axis/.style={-{Stealth[length=4pt]},line width=.65pt,draw=ink!65},
vec/.style={-{Stealth[length=5pt]},line width=1pt,draw=blue},
short/.style={-{Stealth[length=3pt]},line width=.8pt},
guide/.style={dashed,draw=ink!40,line width=.5pt},
card/.style={rounded corners=2pt,draw=blue!65,fill=blue!4,align=center,text width=6.5cm,inner xsep=9pt,inner ysep=7pt},
title/.style={font=\fontsize{10.5}{14}\selectfont\bfseries,text width=15cm,align=center},
label/.style={font=\fontsize{8.4}{11}\selectfont,align=center},
note/.style={label,text width=6.2cm}}
\begin{document}
\begin{tikzpicture}[x=1cm,y=1cm]
\path[use as bounding box] (-8,-3.5) rectangle(8,3.5);
\node[title] at (0,3.15) {M3-030\quad 带自由面电荷的介质分界面};
\begin{scope}[shift={(-4,0)}]
\fill[green!6] (-2.4,0) rectangle(2.4,2.1);
\fill[blue!5] (-2.4,-1.9) rectangle(2.4,0);
\draw[ink!65,thick] (-2.4,0)--(2.4,0);
\node[label] at(-1.8,-1.6){介质1 $\varepsilon_1$};\node[label] at(-1.8,1.8){介质2 $\varepsilon_2$};
\draw[vec,orange] (-1.8,-.4)--(-1.8,.7) node[left] {$\mathbf n$};
\draw[vec,orange] (-.8,0)--(.3,0) node[midway,below] {$\mathbf t$};
\draw[vec] (.9,-1.65)--(.9,-.25) node[midway,right] {$\mathbf D_1$};
\draw[vec] (.9,.25)--(.9,2) node[midway,right] {$\mathbf D_2$};
\draw[red,dashed] (.3,-.45) rectangle(1.5,.45);
\node[label,red] at(1.7,0){$+\sigma_f$};\end{scope}
\node[card,draw=blue!65,fill=blue!4] at (3.9,1.8) {$D_{1n}=6.0\times10^{-9}\ \mathrm{C/m^2}$};
\node[card,draw=green!65,fill=green!4] at (3.9,0) {$D_{2n}=D_{1n}+\sigma_f$\\$=8.0\times10^{-9}\ \mathrm{C/m^2}$};
\node[card,draw=orange!65,fill=orange!4] at (3.9,-1.8) {$\displaystyle E_{1n}=\frac{3.0\times10^{-9}}{\varepsilon_0}$\\$\displaystyle E_{2n}=\frac{1.6\times10^{-9}}{\varepsilon_0}$};
\end{tikzpicture}\end{document}